% !TEX root = ../nc-cubic.tex
\section{点スキームと Poisson 零スキームを比較する際の障害}
\label{ppz:sec:comparison}

点スキームと Poisson tensor の零スキームとの一致を検証する際には、
抽象スキームとしての同型と、既存の次数一偏極に関する埋込み同型を区別する。
まず一ブロック型の零スキームを計算し、埋込み同型に対する障害を示す。
次に長さ八の局所 Artin 環と Hessian 階数二の幾何を用い、
複素数体上の具体的な非可換三次三様体について、
同じ Kuznetsov 成分を持つ全ての可換三次と全ての Poisson tensor を調べる。
その例では可換 Kuznetsov 実現が直接構成できる一方、
全点スキームと零スキームの抽象同型は存在しない。
圏比較には系\ref{csc:cor:fermat}の分離テンソル比較を用いる。

\subsection{射影 Poisson tensor の零イデアル}

$k$ を標数零の代数閉体、$n\geq3$、$r=n-1$ とし、
$V=\langle x,y,z,u_1,\ldots,u_r\rangle$、
$K=\operatorname{ann}\langle u_1,\ldots,u_r\rangle\subset V^*$ と置く。
$H\in\Sym^3V$ に対し、$X=V(H)\subset\PP(V^*)$ と
$g_H=H|_{u=0}$ の定める平面三次がともに滑らかとする。
$0\ne P\in\det K$ を選び、
\[
 \Pi_H=\iota_{dH}P,\qquad \pi_H=\Pi_H|_X
\]
とする。$Z_X(\pi_H)$ は、切断 $\pi_H\in H^0(X,\bigwedge^2T_X)$
の零スキーム、すなわち
$\bigwedge^2\Omega_X\xrightarrow{\pi_H}\OO_X$ の像イデアルで定義する。

\begin{proposition}[一ブロック Poisson 零スキームの分解]
\label{ppz:prop:zero-scheme}
上の仮定の下で、スキームとして
\begin{equation}\label{ppz:eq:zero-decomposition}
 Z_X(\pi_H)=E_H\amalg D_H,\qquad
 E_H=V(u_1,\ldots,u_r,H),\quad
 D_H=V(H_x,H_y,H_z,H)
\end{equation}
となる。
$C_H=V(H_x,H_y,H_z)$ 上で $H$ が非零因子ならば、
$D_H$ は型 $(2,2,2,3)$、次元 $n-3$、次数 $24$ の完全交叉である。
相対的な滑らかさと、この相対 Cartier 条件を仮定した族では、
表示した零イデアルの同定は任意の基底変更と両立する。
\end{proposition}

\begin{proof}
$P$ の非零定数倍は零イデアルを変えないので、
$P=\partial_x\wedge\partial_y\wedge\partial_z$ としてよい。
射影空間の Euler 列は局所的に分裂し、その外積により、
斉次 bivector の射影像が零となるイデアルは
Euler vector field $\operatorname{Eu}$ との積
$\operatorname{Eu}\wedge\Pi_H$ の係数で計算できる。
また $X$ が滑らかなので $T_X\subset T_{\PP(V^*)}|_X$ は部分束であり、
$\bigwedge^2T_X$ の切断を周囲の外積束で計算しても零イデアルは変わらない。
従って $X$ 上の零イデアルは
\begin{equation}\label{ppz:eq:zero-ideal}
 \bigl(u_aH_x,u_aH_y,u_aH_z\ (1\leq a\leq r),
       xH_x+yH_y+zH_z\bigr)
\end{equation}
で与えられる。

ある $u_a$ が可逆な chart では、このイデアルは
$(H_x,H_y,H_z)$ となり、$D_H$ を得る。
$u=0$ の点では $g_H$ が滑らかなので、
$H_x,H_y,H_z$ の少なくとも一つが可逆な近傍を取れる。
その近傍では最初の $r$ 組の方程式が全ての $u_a$ を消す。
最後の方程式は
\[
 xH_x+yH_y+zH_z=3H-\sum_a u_aH_{u_a}
\]
により $X$ 上で冗長であり、零スキームは $E_H$ となる。
$C_H\cap\PP(K)=\varnothing$ は $g_H$ の滑らかさから従うため、
この二つの閉スキームは交わらない。
二種類の chart は零スキーム全体を覆い、式
\eqref{ppz:eq:zero-decomposition}を得る。

三つの active 偏微分は、定理\ref{thm:relative-point-nfold}の
濾過 Koszul 計算と同様に正則列である。
その商は中心多項式環上階数八で有限自由となる。
非零因子 $H$ を加えると完全交叉型・次元・次数が従う。
相対版の局所計算では、幾何ファイバー上で非零な偏微分を可逆にしている。
各 chart のイデアル等式は係数環上の等式なので、
非平坦・非被約のものを含む基底変更後にも成立する。
\end{proof}

\subsection{同じ埋込みを保つ一致の不存在}

量子側には第\ref{sec:points}節と第\ref{sec:poisson-morphism}節の
良い一ブロック中心商 $A=B/(w)$ を取り、
\[
 \operatorname{Pt}(A)=E\amalg D,\qquad
 Q=I(D)_2=\langle F_x,F_y,F_z\rangle,
 \qquad C=V(Q)
\]
とする。
特に $w\notin\Sym^3U$ とし、境界 Jacobi 線への中心評価
\[
 \epsilon:Z(B)_3/\Sym^3U\longrightarrow
       \bigl(k[x,y,z]/(g_x,g_y,g_z)\bigr)_3
\]
が同型であると仮定する。
この写像は中心三次空間全体の同型ではなく、商線束間の同型である。

\begin{corollary}[元の Poisson tensor を保つ切断の障害]
\label{ppz:cor:fixed-tensor}
元の $\Pi=h\Pi_F$ とその support を固定する。
滑らかな古典 Casimir 三次 $H\in kF\oplus\Sym^3U$ に対し、
$D_H=D$ をスキームとして実現することはできない。
従って trace--Hessian の中心同定だけでは、
定常点成分を Poisson 零点の定常成分と同定できない。
\end{corollary}
\begin{proof}
滑らかな $H$ では $H=cF+G(u)$ の係数 $c$ は非零である。
従って $C_H=C$ であり、$D_H=C\cap V(H)$ となる。
命題\ref{prop:point-casimir-restriction-obstruction}により、
この Cartier 切断は $D=C\cap V(w_{\mathrm{ab}})$ と一致しない。
\end{proof}

\begin{proposition}[良い一ブロック型の埋込み一致に対する障害]
\label{ppz:prop:embedded-obstruction}
上の仮定に加え、$g=F|_{u=0}$ が
式\eqref{eq:point-hessian-scalar}の Hessian スカラー条件を満たすとする。
滑らかな一ブロック Poisson 対 $(X_H,\pi_H)$ について、
active 平面三次が滑らかであり、$D_H$ が上の型 $(2,2,2,3)$ の
完全交叉であると仮定する。
このとき、射影線型同型によって
\[
 \operatorname{Pt}(A)\simeq Z_{X_H}(\pi_H)
\]
を埋込みスキームとして実現することはできない。
ここでは標的の Poisson tensor を元の $\Pi$ に固定していない。
\end{proposition}

\begin{proof}
仮にこの射影同型があれば、標的のデータを引き戻して
同じ $\PP(V^*)$ の中で比較する。
$n=3$ では楕円成分は唯一の正次元成分、$n>4$ では唯一の一次元成分である。
$n=4$ では $D,D_H$ は連結な次数 $24$ の完全交叉曲線であり、
次数 $3$ の連結成分 $E,E_H$ と区別できる。
従って射影同型は $E$ を $E_H$ に、$D$ を $D_H$ に移す。
$E$ と $E_H$ の線型包はそれぞれの support 平面なので、
引戻し後の support は元の $K$ と一致する。

完全交叉イデアルの二次部分から
\[
 \langle H_x,H_y,H_z\rangle=I(D_H)_2=I(D)_2=Q
\]
を得る。
命題\ref{prop:point-casimir}の Hessian スカラー条件による積分の一意性から
$H=cF+G(u)$ である。
よって前系の切断障害に反する。
\end{proof}

\begin{remark}[第二予想との論理的な関係]
この命題が排除するのは、良い一ブロック標的について、
既存の次数一埋込みを射影同型まで保つ強い一致である。
点スキームと零スキームの抽象同型、
別の次数一表示や偏極を用いた比較、
複数ブロックまたは退化した Poisson 標的を含む存在問題を
この命題だけで否定したことにはならない。
特に楕円成分の一致と、定常成分を含む全零スキームの一致は別である。
命題\ref{prop:nfm-correction-point-cubic}が与える Hessian 補正の整合は前者であり、
定理\ref{thm:nfm-residual-formal-poisson}の残余比較も
全定常点スキームとの一致を結論していない。
\end{remark}

\subsection{自然な残余 Poisson 対では抽象スキームも異なり得る}

次の結果は、既に構成した特定の残余 Poisson 対が
第二予想\ref{conj:nc-second}を自動的に実現するという推論を排除する。
別の Poisson tensor の存在を排除する結果ではない。
被約な台では消える差を、定常点の局所 Artin 環で検出する。

\begin{proposition}[具体的な分離型 Hesse 対の定常局所環の不一致]
\label{ppz:prop:split-artin-obstruction}
命題\ref{esi:prop:factor}の複素 Sklyanin 因子と指定中心三次式を
$S_\star,\Omega_\star$ と書く。すなわち
\begin{equation}\label{ppz:eq:explicit-pair}
 \begin{split}
 S_\star&=\C\langle x,y,z\rangle/
 (yz-3zy-2x^2,\ zx-3xz-2y^2,\ xy-3yx-2z^2),\\
 \Omega_\star&=\frac{79}{144}(x^3+y^3+z^3)
       +\frac5{96}(xyz+yzx+zxy+yxz+xzy+zyx),\\
 A_\star&=S_\star[u,v]/(\Omega_\star+G),\qquad
 G=(u^3+v^3)/3.
 \end{split}
\end{equation}
$u,v$ は中心変数である。
\[
 g_t=(x^3+y^3+z^3)/3+txyz,\qquad
 \lambda=17/9,\qquad X_\star=V(g_\lambda+G)\subset\PP^4_\C
\]
と置く。
この滑らかな可換三次の自然な一ブロック tensor を
\[
 \pi_\star
 =\left.\iota_{d(g_\lambda+G)}
       (\partial_x\wedge\partial_y\wedge\partial_z)
   \right|_{X_\star}
\]
とする。その非零スカラー倍でも以下の結論は同じである。
このとき
\begin{equation}\label{ppz:eq:split-artin-nonisom}
 \operatorname{Pt}(A_\star)\not\simeq Z_{X_\star}(\pi_\star)
\end{equation}
が抽象スキームとして成立する。
一方、両者の被約部分は、元の $\PP^4_\C$ の中でも一致する。
さらに実際の dg Morita 同値
$\Ku_{\mathrm{coh}}(A_\star)\simeq\Ku(X_\star)$ が存在する。
\end{proposition}

\begin{proof}
$J_t=\C[x,y,z]/(x^2+tyz,y^2+tzx,z^2+txy)$ と置く。
$1+t^3\ne0$ ならば $J_t$ は Hilbert 級数 $(1+s)^3$ の
長さ八の局所 Artin 環であり、極大イデアルの四乗は零である。
三つの非可換二次関係の可換化は、非零定数倍を除いて
$x^2+yz,y^2+zx,z^2+xy$ である。
$J_1$ では $x^3=y^3=z^3=-xyz$ なので、指定した中心三次式の像は
\[
 \eta=[\Omega_\star]
   =\left(-\frac{79}{48}+\frac5{16}\right)xyz
   =-\frac43xyz\ne0.
\]
従って境界中心評価の非消滅条件も満たす。
定理\ref{csp:thm:point-functor}により、点スキームは
$E_\lambda\amalg D_\star$ であり、その定常成分は
$D_\star=\operatorname{Proj}(J_1[u,v]/(G+\eta))$ である。
ここで $E_\lambda=V(g_\lambda)$ は命題\ref{esi:prop:factor}で確認した点曲線である。
定常成分の被約部分は
$x=y=z=0$ と $u^3+v^3=0$ で与えられる三点である。
全ての点は $u\ne0$ に属するので、$u=1$ とする。
各 $a\in\C$、$a^3=-1$ に対して
\[
 v=a-\eta/a^2
\]
は $J_1$ の中で $(1+v^3)/3+\eta=0$ を厳密に満たす。
これは $\eta^2=0$ による。
三つの根の差は単元なので、可換環上の中国剰余定理から
\begin{equation}\label{ppz:eq:source-local-artin}
 D_\star\simeq\coprod_{a^3=-1}\operatorname{Spec}J_1
\end{equation}
を得る。従って中心三次の非零 socle 項を捨てたわけではなく、
その項を中心座標の冪零補正に吸収して局所環を計算した。

命題\ref{ppz:prop:zero-scheme}を $(X_\star,\pi_\star)$ へ適用すると、
その定常成分は
\[
 D_{\mathrm{cl}}=V((g_\lambda)_x,(g_\lambda)_y,(g_\lambda)_z,g_\lambda+G)
 \simeq\coprod_{a^3=-1}\operatorname{Spec}J_\lambda
\]
である。最後の式には Euler 恒等式による $[g_\lambda]=0$ in $J_\lambda$
を用いた。

仮に $J_1\simeq J_\lambda$ が局所 $\C$-代数として成立すれば、
極大イデアルの冪に関する associated graded 環も同型になる。
両環は標準次数付きで斉次二次イデアルによる商なので、
この associated graded 環はそれぞれ $J_1,J_\lambda$ 自身である。
次数一と二の乗法から二つの勾配空間の線型同型を得て、
命題\ref{prop:ej-projective-uniqueness}により
$V(g_1)\simeq V(g_\lambda)$ を得るはずである。
しかし Hesse の不変量の公式は厳密に
\begin{equation}\label{ppz:eq:explicit-j}
 \begin{split}
 j(t)&=\frac{27t^3(8-t^3)^3}{(1+t^3)^3},\qquad j(1)=\frac{9261}{8},\\
 j(17/9)&=\frac{3813232609367}{4849120870776},\qquad
 j(17/9)-j(1)=-\frac{1402412578855675}{1212280217694}\ne0
 \end{split}
\end{equation}
を与える。従ってこの局所代数同型は存在しない。
点スキームの一次元成分は楕円曲線一つだけなので、
全スキームの抽象同型があれば零次元成分同士が対応する。
三つ全ての局所環が上の非同型な二種類であるため、
式\eqref{ppz:eq:split-artin-nonisom}が従う。

両側の楕円成分は同じ埋込み平面三次 $E_\lambda$ であり、
定常成分の被約部分も上の同じ三点である。
これで被約部分の一致も得る。
最後に、命題\ref{esi:prop:factor}は実楕円因子の正則性と中心商
$S_\star/(\Omega_\star)=B(E_\lambda,\OO_{E_\lambda}(1),t_p)$ を与える。
系\ref{csc:cor:fermat}を $a=1$、$c=d=1/3$、$f=g_\lambda$ に適用すると、
複素数体上で直接 $\Ku_{\mathrm{coh}}(A_\star)\simeq\Ku(X_\star)$ を得る。
\end{proof}

\begin{remark}[自然な処方への反例と存在問題の区別]
この例は、自然な残余 support と tensor をそのまま使う処方に対する反例である。
$X_\star$ 上で別の support を持つ Poisson tensor を選ぶ可能性や、
同じ Kuznetsov 成分を持つ別の可換三次を選ぶ可能性は扱っていない。
従ってこの命題だけでは「ある可換三次とある Poisson tensor が同時に存在する」という
第二予想\ref{conj:nc-second}を否定したことにはならない。
次小節で任意の support を調べ、この例についてその可能性も排除する。
一方、被約な台や楕円成分だけの比較では失われる情報が、
この例では長さ八の局所環に残っている。
一般パラメータにおける Hesse の有理恒等式は
第\ref{sec:nfm-residual-th-obstruction}節に記載した。
ここで用いた有理数の厳密検算は
\texttt{computations/verify\_point\_poisson\_obstruction.py}に含まれる。
局所環の分解と非同型の証明は上の議論であり、
その有限多項式検算と同一視しない。
\end{remark}

\subsection{三次三様体における任意の Poisson tensor の排除}

三次三様体では、全ての非零 Poisson tensor が分解可能な定数三ベクトルから来る。
これは Mayanskiy\cite[Theorem~1, \S2.2]{Mayanskiy}の既知結果である。
同論文の Schouten 括弧の式では、Pl\"ucker 関係式が五つの偏微分の係数となる。
滑らかな三次の偏微分は一次独立なので、Poisson 条件は全 Pl\"ucker 関係に等しい。
従ってこの結論は係数の特殊な三次にも適用でき、
同じ多項式計算によって標数零の代数閉体上でも成立する。
以下ではこの分類と、孤立零点の局所環を組み合わせる。

\begin{lemma}[埋込み次元三の孤立 Poisson 零点]
\label{ppz:lem:isolated-artin-eckardt}
$k$ を標数零の代数閉体、$H\in k[W]_3$、$\dim W=5$ とし、
$X=V(H)\subset\PP(W)$ を滑らかな三次三様体とする。
$0\ne\pi\in H^0(X,\bigwedge^2T_X)$ を Poisson tensor とし、
$p\in Z_X(\pi)$ を、局所環が長さ八・埋込み次元三を持つ孤立点とする。
$p$ の非零ベクトル代表を同じ文字で表す。
このとき $B_p=D^2H(p)$ の階数は二である。
さらに
\[
 R_p=\ker B_p,\qquad G_p=H|_{R_p}
\]
と置くと、$\dim R_p=3$、$V(G_p)\subset\PP(R_p)$ は滑らかな平面三次であり、
局所環の極大イデアル $\mathfrak m$ に対して
\begin{equation}\label{ppz:eq:eckardt-graded-local}
 \operatorname{gr}_{\mathfrak m}\OO_{Z_X(\pi),p}
 \simeq\operatorname{Jac}(G_p)
\end{equation}
となる。右辺は定数倍および線型座標変換を除いて $X,p$ だけから定まる。
\end{lemma}

\begin{proof}
上の分類により、$\pi=(\iota_{dH}P)|_X$、
$0\ne P\in\det K$、$\dim K=3$ と書ける。
$X\cap\PP(K)$ は Poisson 零集合に含まれる。
実際、この平面上では Euler vector と bivector の外積は
$3H P$ であり、$X$ 上で零になる。
この交叉は平面曲線、または平面全体なので、
孤立点 $p$ は $\PP(K)$ に属さない。

$p$ で非零な中心一次式を一にした affine chart を取る。
零点条件は $H_x=H_y=H_z=0$ であるから、
三つの $K$ 方向は $p$ で $X$ の接方向である。
滑らかさにより残りの affine 座標 $s$ を $H_s(p)\ne0$ となるよう選べ、
$x,y,z$ は $X$ の局所座標となる。
零イデアルはその正則局所環の中で $(H_x,H_y,H_z)$ である。
商の埋込み次元が三であることは、これらの元に一次項がないことを意味する。
従って $B_p|_{K\times K}=0$ である。

affine cone の接超平面を $T=\ker dH(p)\subset W$ とすると、
$T=K\oplus kp$ である。
Euler の式から $B_p(p,v)=2dH(p)(v)$、$B_p(p,p)=6H(p)=0$ なので、
$B_p|_{T\times T}=0$ となる。
余次元一の部分空間上で零となる対称形式の階数は高々二であり、
実際には二である。階数一なら $B_p=c\ell^2$ と書け、
$B_p(p,p)=0$ から $\ell(p)=0$、従って
$B_p(p,-)=0$ となって $dH(p)\ne0$ に反する。
従って $R_p$ は三次元で、$R_p\subset T$、$p\notin R_p$、
\[
 T=R_p\oplus kp
\]
を得る。

$B_p|_{T\times T}=0$ と $H(p)=0$ により、$H|_T$ は $p$ 方向に依らず
$H(v+tp)=H(v)$ for $v\in R_p$ である。
$G_p$ が特異と仮定し、$0\ne a\in R_p$ をその特異点の代表とする。
全ての $t$ について $a+tp\in V(H)$ であり、
その点での $T$ 方向の微分は全て零である。
$n\notin T$ を取ると、残る微分
$dH(a+tp)(n)$ は $t$ の二次式で、その最高次係数は
$dH(p)(n)\ne0$ である。
代数閉体上でその零点を選ぶと $X$ の特異点となり、矛盾する。
よって $G_p$ は滑らかである。

$K$ は $T=R_p\oplus kp$ の $kp$ に対する補空間だから、
$R_p$ の $p$ 方向への graph であり、$H|_K$ は $G_p$ と線型同型である。
$p$ を原点とする上の affine chart では、方程式を
\[
 H(x,y,z,s)=G_p(x,y,z)+sA(x,y,z,s),\qquad A(0)\ne0
\]
と書ける。形式的に $s$ を解くと $s\in(x,y,z)^3$ であり、
三つの零イデアルの生成元の二次先頭項は
$(G_p)_x,(G_p)_y,(G_p)_z$ である。
従って $\operatorname{Jac}(G_p)$ から
$\operatorname{gr}_{\mathfrak m}\OO_{Z_X(\pi),p}$ への次数付き全射がある。
前者は滑らかな平面三次の勾配正則列による長さ八の環であり、
後者も仮定によって長さ八である。
従ってこの全射は同型で、式\eqref{ppz:eq:eckardt-graded-local}を得る。
\end{proof}

\begin{lemma}[分離型 Hesse 三様体の階数二 Hessian 点]
\label{ppz:lem:hesse-eckardt}
$\lambda\ne0$、$\lambda^3\ne8,-1$ とし、
\[
 H_\lambda=g_\lambda+(u^3+v^3)/3
\]
と置く。
$V(H_\lambda)$ 上で Hessian 行列の階数が高々二となる点は、
中心直線上の三点と active 平面上の九つの flex、合計十二点だけである。
前補題の平面三次 $G_p$ の $j$ 不変量は、
中心三点では $j(g_\lambda)$、九つの flex では $0$ である。
\end{lemma}

\begin{proof}
Hessian 行列は
\[
 \operatorname{diag}\bigl(M_\lambda(x,y,z),2u,2v\bigr),\qquad
 M_\lambda=\begin{pmatrix}
 2x&\lambda z&\lambda y\\
 \lambda z&2y&\lambda x\\
 \lambda y&\lambda x&2z
 \end{pmatrix}
\]
である。$(x,y,z)\ne0$ では $M_\lambda$ の階数は二以上である。
実際、階数が一以下で一座標が零なら、主小行列式と $\lambda\ne0$ により
残りの二座標も零となる。
全て非零なら、非主小行列式から $2x^2=\lambda yz$、
主小行列式から $4yz=\lambda^2x^2$ となり、$\lambda^3=8$ に反する。

従って active 座標が非零な階数二点では $u=v=0$ である。
$g_\lambda=0$ 上では
\[
 \det M_\lambda=8(1+\lambda^3)xyz
\]
なので、該当点は一座標が零となる九点、すなわち九つの flex である。
active 座標が零なら $u^3+v^3=0$ であり、
$u,v$ はともに非零で、中心直線上の三点を得る。

中心三点では Hessian の核は active 三平面であり、$G_p=g_\lambda$ である。
flex は座標の置換により
$p=(1,-\epsilon,0,0,0)$、$\epsilon^3=1$ としてよい。
このとき Hessian の核は $e_u,e_v$ と
\[
 r=(\lambda\epsilon/2,\lambda\epsilon^2/2,1,0,0)
\]
で張られ、その制限は
\[
 H_\lambda(sr+u e_u+v e_v)
 =\bigl((1+\lambda^3)s^3+u^3+v^3\bigr)/3
\]
となる。これは滑らかな Fermat 平面三次と同型で、$j=0$ である。
\end{proof}

\begin{theorem}[分離型三次三様体における全 Poisson 零スキームの障害]
\label{ppz:thm:all-poisson-obstruction}
$k$ を標数零の代数閉体とし、
\[
 \lambda\ne0,\qquad \lambda^3\ne8,-1,\qquad
 1+t^3\ne0,\qquad j(t)\notin\{0,j(\lambda)\}
\]
とする。
$X_\lambda=V(g_\lambda+(u^3+v^3)/3)$ 上のどの Poisson tensor の
零スキームも、局所環が
$J_t=k[x,y,z]/(x^2+tyz,y^2+tzx,z^2+txy)$ と同型な孤立点を持たない。
従って、そのような孤立点を持つ任意のスキームは、
どの $Z_{X_\lambda}(\pi)$ とも同型でない。

特に命題\ref{ppz:prop:split-artin-obstruction}の具体的な複素代数について、
$X_\star$ 上のどの Poisson tensor $\pi$ に対しても
\[
 \operatorname{Pt}(A_\star)\not\simeq Z_{X_\star}(\pi)
\]
である。ここでは support の指定、active 制限の滑らかさ、
零スキームの完全交叉性を標的への追加仮定として課さない。
\end{theorem}

\begin{proof}
$\pi=0$ の零スキームは滑らかな三次元の $X_\lambda$ 全体であり、孤立点を持たない。
$\pi\ne0$ の場合、仮に $J_t$ と同型な局所環を持つ孤立零点があれば、
その長さは八、埋込み次元は三である。
補題\ref{ppz:lem:isolated-artin-eckardt}と
補題\ref{ppz:lem:hesse-eckardt}により、その associated graded 環は
$J_\lambda$ または Fermat 三次の Jacobian 環になる。
$J_t$ は標準次数付きなので、自身の極大イデアル filtration の
associated graded 環に等しい。
よって命題\ref{prop:ej-projective-uniqueness}から
$j(t)=j(\lambda)$ または $j(t)=0$ が必要になり、仮定に反する。

具体例では $t=1$、$\lambda=17/9$ である。
式\eqref{ppz:eq:explicit-j}に加えて
\[
 \lambda^3-8=-\frac{919}{729}\ne0,\qquad
 \lambda^3+1=\frac{5642}{729}\ne0
\]
だから全ての仮定を満たす。
式\eqref{ppz:eq:source-local-artin}の三つの孤立点により、最後の結論を得る。
\end{proof}

この判定は固定された有理パラメータの例だけに限らない。
第\ref{sec:nfm-residual-th-obstruction}節の Hesse 有理恒等式を用いれば、
同節の幾何学的一般点でも同じ $j$ の非一致条件が成立する。
この一般パラメータでの計算と、以下の複素数体上の直接構成を区別する。

\subsection{具体的な複素非可換三次三様体による反例}

\begin{corollary}[全零スキーム版の第二予想に対する複素数体上の反例]
\label{ppz:cor:complex-counterexample}
式\eqref{ppz:eq:explicit-pair}の $A_\star$ は、複素数体上の
点スキームの良い一ブロック条件を満たす滑らかな楕円一ブロック型の
非可換三次三様体であり、次の二つを同時に満たす。
\begin{enumerate}
 \item 滑らかな可換三次三様体
 \[
 X_\star=V\!\left(
 \frac{x^3+y^3+z^3+u^3+v^3}{3}+\frac{17}{9}xyz\right)
 \subset\PP^4_\C
 \]
 と実際の dg Morita 同値
 $\Ku_{\mathrm{coh}}(A_\star)\simeq\Ku(X_\star)$ が存在する。
 \item $\Ku_{\mathrm{coh}}(A_\star)\simeq\Ku(Y)$ を満たす
 どの滑らかな可換三次三様体 $Y$ と、
 その上のどの Poisson tensor $\pi$ に対しても、
 $\operatorname{Pt}(A_\star)\not\simeq Z_Y(\pi)$ が抽象スキームとして成立する。
\end{enumerate}
従って、非被約構造を含む全 point scheme と全 Poisson 零スキームの同型を要求する
第二予想\ref{conj:nc-second}は成立しない。
主予想の可換 Kuznetsov 実現はこの例では成立している。
\end{corollary}

\begin{proof}
系\ref{esi:cor:ambient}は、指定した複素代数の AS 正則な周囲環と
中心正則三次商を与える。
命題\ref{ppz:prop:split-artin-obstruction}の証明で、
系\ref{csc:cor:fermat}による全 coherent 残余圏の比較を構成した。
従って第一項が成立し、同じ比較の smoothness と coherent 飽和により
$A_\star$ は滑らかな非可換三次三様体である。
active 可換化 $g_1$、点曲線 $g_{17/9}$ はともに滑らかであり、
中心三次の socle 評価 $-4xyz/3$ は非零である。
定常三次は $J_1[u,v]$ 上の非零因子であるから、
点スキームの良い一ブロック条件も満たす。
ここでいう良い条件は点関手の条件であり、
全 $\OO$-表示 stack の全ての相対入力への所属を同時に主張するものではない。
それらの入力は本系の反例の証明には用いない。

$X_\star$ 上の全ての Poisson tensor は
定理\ref{ppz:thm:all-poisson-obstruction}により排除される。
別の標的 $Y$ に対しては、第一項の同値を合成して
$\Ku(Y)\simeq\Ku(X_\star)$ を得た後、
三次三様体の categorical Torelli
\cite[Theorem~1.1]{NFMThreefoldCategoricalTorelli}から $Y\simeq X_\star$ を得る。
$Y$ 上の Poisson tensor をこの同型で $X_\star$ へ移せば、
その零スキームも移送されるため、第二項が従う。
\end{proof}

\begin{remark}[反例が排除する範囲と残る問題]
この反例は $n=3$ の非被約な定常点を持つ分離型であり、
被約 point scheme を持つ一般混合族の一致や、被約化後の一致を否定しない。
実際、命題\ref{ppz:prop:split-artin-obstruction}では
自然な残余 Poisson 対と被約部分が一致する。
従って被約化を取った命題、滑らかな定常成分を仮定した命題、
または別の点関手による命題は、別の検証問題として残る。
これらのいずれかを元の第二予想の代わりに自動採用するものではない。
Hessian と制限三次の恒等式は
\texttt{computations/verify\_point\_poisson\_obstruction.py}で検算する。
全 support の排除、局所環、Ku 比較、Torelli の適用は本文で証明し、
この計算で代用しない。
\end{remark}

